\teachcover{Changing the arithmetic}{1}{
Rings, fields, and vectors
\[
X^2-X-1=0\qquad\text{modulo }p
\]
}{
The students already know Fibonacci periodicity via a matrix. Today names
the arithmetic they have been using and begins the vector-space reminder.
The goal is not diagonalization yet. Allow calculation time at slides 8
and 20. All rings in these lectures are commutative with identity, with
$1\ne0$. A field has the additional inverse property.
}

\teachslide{The matrix we already know}{2}{
\[
A=\mat1110,\qquad
\begin{pmatrix}F_{n+1}\\F_n\end{pmatrix}
=A^n\begin{pmatrix}1\\0\end{pmatrix}.
\]
\[
\det A=-1,\qquad A^2=A+I.
\]
Modulo any integer $m\ge2$, the step is reversible.

\key{Today: what arithmetic are these entries using?}
}{
Use $F_0=0,F_1=1$. This matrix orders the state as $(F_{n+1},F_n)$;
the book sometimes uses the reversed state and the conjugate matrix.
They give the same period. Ask students to multiply $A$ by itself.
The determinant is a unit modulo every modulus, so the inverse exists.
Do not confuse invertibility of this particular matrix with every
nonzero residue being invertible.
}

\teachslide{Where are the roots?}{2}{
\[
\det(XI-A)=X^2-X-1.
\]
Over the real numbers, its roots are
\[
\frac{1+\sqrt5}{2},\qquad \frac{1-\sqrt5}{2}.
\]
\key{Modulo $11$:}\quad $4^2=5$.

\key{Modulo $7$:}\quad squares are $0,1,2,4$.
}{
Reduce $16$ modulo $11$. For $7$, list the squares of $0,1,2,3$;
the remaining inputs give the same squares. Stress that writing a
square-root symbol does not manufacture a residue with that property.
Division by $2$ will make sense for odd primes. Prime $2$ is exceptional
and can wait until Lecture 3's optional material.
}

\teachslide{Residues are numbers with arithmetic}{2}{
\[
\Z/n\Z=\{[0],[1],\ldots,[n-1]\}.
\]
\[
[a]+[b]=[a+b],\qquad [a][b]=[ab].
\]
For example, in $\Z/6\Z$,
\[
[4]+[5]=[3],\qquad [4][5]=[2].
\]
We usually omit the brackets once the modulus is fixed.
}{
The objects are residue classes, not just chosen representatives.
Recall that replacing $a$ by $a+kn$ changes neither the sum class nor
the product class. This is the well-definedness behind modular arithmetic.
Ask for the two reductions on the slide before displaying the next page.
}

\teachslide{The ring rules}{3}{
A \key{commutative ring with identity} has addition and multiplication.
\begin{itemize}
\item Addition is associative and commutative; it has $0$ and negatives.
\item Multiplication is associative and commutative; it has $1\ne0$.
\item Multiplication distributes over addition.
\end{itemize}
\[
a(b+c)=ab+ac.
\]
}{
State closure under both operations. Say explicitly that $a+0=a$,
$a+(-a)=0$, and $a\cdot1=a$. Associativity allows us to omit parentheses.
Integers and residue classes are the running examples; real numbers also
form a ring. These are the commutative rings with identity used in the
course; other conventions exist. Do not add the field inverse axiom yet.
The calculations inherited from integer arithmetic justify the axioms
for $\Z/n\Z$.
}

\teachslide{A new phenomenon: zero divisors}{2}{
In $\Z/6\Z$,
\[
\key{2\cdot3=0,\qquad 2\ne0,\quad3\ne0.}
\]
A \key{zero divisor} is a nonzero element $a$ for which
\[
ab=0\quad\text{for some nonzero }b.
\]
Cancellation can fail:
\[
2\cdot1=2\cdot4,\qquad 1\ne4.
\]
}{
Check both products in the last equation: $2$ and $8$ have the same
residue modulo $6$. Link cancellation failure to $2(1-4)=0$.
Ask whether this could happen in the rational numbers. The point is
that the arithmetic laws remain valid while familiar cancellation
requires extra hypotheses.
}

\teachslide{Which elements can we divide by?}{2}{
A \key{unit} is an element $u$ with a multiplicative inverse:
\[
uv=1\quad\text{for some }v.
\]
In $\Z/n\Z$,
\[
\key{[a]\text{ is a unit}\iff\gcd(a,n)=1.}
\]
In $\Z/6\Z$, the units are $1$ and $5$.
}{
Recall Bezout: $ax+ny=1$ supplies the inverse $[x]$. Conversely,
$ax\equiv1\pmod n$ gives an integer equation $ax+ny=1$, forcing the gcd
to be $1$. Thus inverses are not a new computational trick: extended
Euclid already computes them. A unit cannot be a zero divisor, since
multiplication by its inverse would force $b=0$.
}

\teachslide{Your turn: modulo eight}{3}{
Work in $\Z/8\Z$.
\begin{enumerate}
\item Which of $1,2,\ldots,7$ are units?
\item Find the inverse of $3$.
\item Find two nonzero elements whose product is $0$.
\end{enumerate}
\spacer
\key{Take one minute; then compare answers.}
}{
Keep this page visible while students calculate. Give a minute of
quiet work, then invite answers in chat or aloud. If needed, suggest
using the gcd test rather than trying every product. Ask students to
reduce $3^2$ and $2\cdot4$. The following page contains the answers.
}

\teachslide{Modulo eight: answers}{1}{
\[
\text{Units: }1,3,5,7.
\]
\[
3^{-1}=3,\qquad 3\cdot3=9=1\pmod8.
\]
\[
2\cdot4=0\pmod8.
\]
\key{Being nonzero does not guarantee an inverse.}
}{
Confirm that every odd residue is coprime to $8$. The nonzero even
residues are zero divisors. This is also a useful reminder that an
element may be its own inverse without equaling $1$.
}

\teachslide{A field allows division}{2}{
A \key{field} is a commutative ring with identity in which
\[
\text{every nonzero element has a multiplicative inverse.}
\]
Examples: $\Q$, $\R$, and $\Z/7\Z$.

A field has no zero divisors:
\[
ab=0,\ a\ne0\quad\Longrightarrow\quad b=a^{-1}0=0.
\]
}{
Division by $a\ne0$ means multiplication by $a^{-1}$. Zero has no
inverse because $0b=0\ne1$. Integers are a ring but are not a field;
for example $2$ has no integer inverse. Do not claim the converse
that every ring without zero divisors is a field: the integers refute it.
}

\teachslide{Why prime moduli give fields}{3}{
\[
\key{\Z/n\Z\text{ is a field}\iff n\text{ is prime}.}
\]
If $n=p$ is prime, then $1\le a<p$ implies $\gcd(a,p)=1$.

So every nonzero class has an inverse.
\spacer
If $n=rs$ with $1<r,s<n$, then
\[
[r][s]=[0]\qquad\text{with }[r],[s]\ne[0].
\]
We write $\F_p=\Z/p\Z$.
}{
Give both directions. In the composite case, $r$ cannot be invertible,
since an inverse would turn $rs=0$ into $s=0$. The statement assumes
$n\ge2$. Use $p=7$ and $n=8$ as the examples students have just seen.
Introduce the blackboard-bold F notation slowly: this is the same
familiar residue arithmetic, now recognized as a field.
}

\teachslide{Characteristic: count copies of one}{2}{
Write $m\cdot1$ for the sum of $m$ copies of $1$.

The \key{characteristic} is the least positive $m$ such that
\[
m\cdot1=0.
\]
If no such $m$ exists, the characteristic is $0$.
\[
\operatorname{char}(\Q)=0,\qquad
\operatorname{char}(\Z/n\Z)=n.
\]
}{
Distinguish repeated addition from an external scalar multiplication
that has not yet been defined. Characteristic zero means that no
positive number of copies of $1$ is zero, not that zero copies sum to
zero. In residues the least positive return is exactly the modulus.
Characteristic is defined for our rings as well as fields.
}

\teachslide{A field's positive characteristic is prime}{3}{
Suppose the characteristic of a field is $m=rs$, with $1<r,s<m$.
\[
(r\cdot1)(s\cdot1)=m\cdot1=0.
\]
Minimality of $m$ gives
\[
r\cdot1\ne0,\qquad s\cdot1\ne0.
\]
This contradicts the absence of zero divisors.
\spacer
\key{A field has characteristic $0$ or a prime.}
}{
Use distributivity to explain why the product of the two repeated sums
is the repeated sum with $rs$ terms. The minimality step is essential:
neither factor is zero since both indices are smaller than $m$.
Remind students that $m=1$ is excluded by $1\ne0$. The ring
$\Z/6\Z$ shows why the field hypothesis is needed.
}

\teachslide{A finite field cannot have characteristic zero}{2}{
In a finite field $K$, the list
\[
0,\ 1,\ 2\cdot1,\ 3\cdot1,\ldots
\]
must contain a repetition.

If $r\cdot1=s\cdot1$ with $r<s$, subtract:
\[
(s-r)\cdot1=0.
\]
\key{Therefore $\operatorname{char}(K)=p$ for some prime $p$.}
}{
This uses finiteness and subtraction, then the preceding theorem.
Do not assume the field's number of elements equals its characteristic.
That tempting assumption is exactly what Lecture 2 will disprove.
Here we know only that the characteristic is some positive prime.
}

\teachslide{Every finite field contains a prime field}{2}{
If $\operatorname{char}(K)=p$, the elements
\[
0,\ 1,\ 2\cdot1,\ldots,(p-1)\cdot1
\]
form a copy of $\F_p$ inside $K$.

Addition and multiplication agree with arithmetic modulo $p$.
\spacer
\key{$K$ may have more elements than these $p$ elements.}
}{
The map $[a]\mapsto a\cdot1$ is well defined modulo $p$, respects the
two operations, and is injective by minimality of the characteristic.
The nonzero elements have inverses within this copy by the same Bezout
argument as before. This copy is called the prime subfield. Prepare
the next step: elements of this smaller field can act as scalars on $K$.
}

\teachslide{Vectors and scalars}{2}{
A \key{vector space over a field $k$} has
\begin{itemize}
\item vectors $u,v,\ldots$ that we can add;
\item scalars $a,b,\ldots\in k$ that multiply vectors.
\end{itemize}
\[
v\longmapsto av.
\]
Examples: $\R^2$ over $\R$ and $\F_3^2$ over $\F_3$.
}{
Always name both the vector space and its scalar field. Vectors need
not be arrows and scalars need not be real numbers. For now use pairs
to keep every operation visible. Later the vectors will be elements
of a larger field. Scalar multiplication is an operation with one
scalar input and one vector input, rather than two vector inputs.
}

\teachslide{The vector-space rules}{3}{
Vector addition has the usual associative and commutative laws,
a zero vector, and additive inverses.

For $a,b\in k$ and $u,v\in V$,
\[
\begin{aligned}
a(u+v)&=au+av, & (a+b)v&=av+bv,\\
a(bv)&=(ab)v, & 1v&=v.
\end{aligned}
\]
\key{The field $k$ supplies the scalar arithmetic.}
}{
State closure under vector addition and scalar multiplication.
Read each law aloud and say which symbols are scalars and which are
vectors. The zero vector is distinct in type from the scalar $0$,
even when we use the same symbol. In coordinate spaces these axioms
hold because the field laws hold in every coordinate. This is a
definition, not an assertion that vectors can be multiplied together.
}

\teachslide{A familiar calculation over a different field}{2}{
In $\F_3^2$, let $u=(1,2)$ and $v=(2,2)$.
\[
u+v=(0,1),\qquad 2u=(2,1).
\]
\[
2u+v=(1,0).
\]
All coordinates are reduced modulo $3$.
\spacer
\key{Same vector rules; different scalar arithmetic.}
}{
Calculate each coordinate before reducing: $(3,4)$, $(2,4)$, and
$(4,3)$. Ask students to verify $2(u+v)=2u+2v=(0,2)$.
Three copies of any vector in this example sum to the zero vector.
Contrast that with $\R^2$ without relying on a geometric picture.
}

\teachslide{Linear combinations and span}{3}{
A \key{linear combination} of $v_1,\ldots,v_r$ is
\[
a_1v_1+\cdots+a_rv_r,\qquad a_i\in k.
\]
Their \key{span} is the set of all such combinations.

In $\F_3^2$, for $u=(1,2)$,
\[
\Span(u)=\{(0,0),(1,2),(2,1)\}.
\]
}{
There are three scalar choices here. The span includes zero and the
original vector, and is closed under addition and scalar multiplication.
For pairs $e_1=(1,0)$, $e_2=(0,1)$, every $(a,b)$ equals $ae_1+be_2$.
Mention this as a contrast: one vector need not span the whole space.
Independence and uniqueness of coordinates are Lecture 2's task.
}

\teachslide{Your turn: a span over $\F_2$}{3}{
In $\F_2^2$, let $u=(1,1)$ and $v=(1,0)$.
\begin{enumerate}
\item List the vectors in $\Span(u)$.
\item List all vectors $au+bv$ with $a,b\in\F_2$.
\item Can $(0,1)$ be written as $au+bv$?
\end{enumerate}
\spacer
\key{Remember: the only scalars are $0$ and $1$.}
}{
Give at least a minute of quiet calculation. Encourage a four-row
table for $(a,b)=(0,0),(1,0),(0,1),(1,1)$. For the third question,
ask for coefficients, not merely a yes. This exercise prepares the
coordinate-counting argument without introducing independence early.
}

\teachslide{Four coefficient choices, four vectors}{1}{
\[
\begin{array}{c|c}
(a,b)&au+bv\\\hline
(0,0)&(0,0)\\
(1,0)&(1,1)\\
(0,1)&(1,0)\\
(1,1)&(0,1)
\end{array}
\]
\[
\Span(u)=\{(0,0),(1,1)\},\qquad (0,1)=u+v.
\]
}{
The table shows both spanning and uniqueness in this example.
We will name the combination of those properties a basis next time.
Do not generalize distinct coefficient choices giving distinct
vectors to an arbitrary spanning list: independence is needed.
}

\teachslide{What we have built today}{2}{
\begin{itemize}
\item Residue arithmetic gives rings; prime moduli give fields.
\item Every finite field contains some $\F_p$.
\item Vector spaces can use $\F_p$ as their scalar field.
\end{itemize}
\spacer
Next time: \key{bases, counting, and a field with $49$ elements.}
}{
Ask students to explain the distinction between a unit and a nonzero
element. Then connect the three bullets: the copy of $\F_p$ supplies
scalars for a vector space whose vectors are elements of the larger
field. We will check its vector-space rules next time rather than
asking them to accept the slogan on faith.
}

\teachslide{Exit question}{2}{
Why is $\Z/6\Z$ a ring but not a field?

Why can it not serve as the scalar \emph{field} in our definition
of a vector space?
\spacer
\key{Give a specific arithmetic obstruction.}
}{
Expected answer: $2\cdot3=0$ with both factors nonzero, so $2$ and $3$
cannot have inverses. Our vector-space definition requires a field
of scalars. Addition and scalar multiplication over a ring can still
be studied; those objects are modules, an optional observation on the
backup slide. End here at minute 50.
}

\backups
\backupslide{Rings as scalars: a glimpse beyond vector spaces}{
One can use a ring of scalars and keep the vector-space identities.

The resulting object is called a \key{module}.

Over $\Z/6\Z$, $2v=0$ need not imply $v=0$:
\[
v=3,\qquad 2v=0.
\]
Our next lectures use fields, where nonzero scalars are invertible.
}{
Use only if a student asks whether the ring case is meaningless.
It is meaningful, but some basis and dimension facts change. The
example uses the ring itself as a module over itself. Avoid developing
module theory during the main lecture.
}

\backupslide{One set, two different scalar fields}{
The complex numbers form a vector space over $\R$:
\[
z=a+bi,\qquad a,b\in\R.
\]
The same set forms a vector space over $\mathbb C$:
\[
z=z\cdot1.
\]
\key{Changing the scalars can change the dimension.}
}{
This example is optional because complex numbers may not be familiar.
Over the reals, $1,i$ form a basis; over the complex numbers, $1$ alone
does. Lecture 2 supplies the precise basis definition. The example
anticipates $\F_{49}$ over $\F_7$ without implying any embedding into
the ordinary real or complex fields.
}
